\documentclass[12pt,a4paper]{article}
\usepackage[utf8]{inputenc}
\usepackage[spanish]{babel}
\usepackage[top=2.5cm, bottom=2.5cm, left=3cm, right=3cm]{geometry}
\usepackage{setspace}
\onehalfspacing
\usepackage{graphicx}
\usepackage{fancyhdr}
\usepackage{enumitem}
\usepackage{booktabs}
\usepackage{hyperref}
\usepackage{tabularx}
\usepackage{amssymb}

\setlength{\headheight}{14.5pt}

\pagestyle{fancy}
\fancyhf{}
\lhead{\textit{Redes de Alta Velocidad}}
\rhead{\textit{Actividad 1 -- Reporte de Lectura}}
\cfoot{\thepage}
\renewcommand{\headrulewidth}{0.4pt}

\makeatletter
\renewcommand{\@seccntformat}[1]{\csname the#1\endcsname.\quad}
\makeatother

\setlist[itemize]{label=\footnotesize$\blacksquare$, itemsep=0.15\baselineskip, topsep=0.5\baselineskip}
\setlist[enumerate]{label=\alph*), itemsep=0.15\baselineskip, topsep=0.5\baselineskip}

\begin{document}

% ==========================================
% PORTADA
% ==========================================
\begin{titlepage}
    \centering
    \vspace*{1cm}
    {\Large\bfseries [Nombre de la Institución]\par}
    \vspace{0.4cm}
    {\large\bfseries [Facultad / Departamento]\par}
    \vspace{1.5cm}
    \noindent\rule{\textwidth}{0.8pt}
    \vspace{1cm}
    {\huge\bfseries Reporte de Lectura\par}
    \vspace{0.7cm}
    {\Large ``Redes de comunicaciones de alta velocidad: \\ ¿cómo funcionan?''\par}
    \vspace{0.6cm}
    {\normalsize José F. Duato Marín. \textit{Rev. R. Acad. Cienc. Exact. Fís. Nat.}, Vol.~109, N.º~1--2, pp.~75--86, 2016.\par}
    \vspace{1cm}
    \noindent\rule{\textwidth}{0.8pt}
    \vspace{1.8cm}

    \begin{tabular}{r l}
        \textbf{Materia:}    & Redes de Alta Velocidad \\[0.35cm]
        \textbf{Actividad:}  & Actividad de apertura -- Lectura 1 \\[0.35cm]
        \textbf{Alumno(a):}  & \underline{\hspace{6.5cm}} \\[0.35cm]
        \textbf{Matrícula:}  & \underline{\hspace{6.5cm}} \\[0.35cm]
        \textbf{Docente:}    & \underline{\hspace{6.5cm}} \\[0.35cm]
        \textbf{Fecha:}      & \underline{\hspace{6.5cm}}
    \end{tabular}

    \vfill
    Ciclo escolar: [XXXX--X]
    \vspace*{0.5cm}
\end{titlepage}
\setcounter{page}{1}

% ==========================================
% ENUNCIADO DE LA ACTIVIDAD
% ==========================================
\section*{Enunciado de la actividad}

Como punto de partida de la materia \textbf{Redes de Alta Velocidad}, se asigna la lectura del artículo de divulgación \textit{``Redes de comunicaciones de alta velocidad: ¿cómo funcionan?''}, de José F. Duato Marín (Rev. R. Acad. Cienc. Exact. Fís. Nat., Vol.~109, 2016). En este texto se presenta, de manera accesible y estructurada, el panorama completo de las interconexiones que sostienen la era de Internet: desde el enlace punto a punto entre dos dispositivos, pasando por la arquitectura interna de un conmutador, hasta las redes multietapa y las redes en chip que integran centenares de núcleos de procesamiento. A lo largo de la lectura se abordan conceptos fundamentales ---control de flujo, gestión de colas, técnicas de conmutación, enrutamiento, interbloqueo, equilibrio de carga, calidad de servicio (QoS) y tolerancia a fallos--- que constituirán el eje vertebrador de los temas que desarrollaremos durante el curso.

Se solicita al estudiante realizar una \textbf{lectura comprensiva} del artículo y elaborar un reporte escrito siguiendo la estructura que se indica a continuación. En particular, deberá \textbf{identificar al menos tres problemáticas de diseño} que surgen al escalar una red (por ejemplo, congestión, interbloqueo y entrega fuera de orden) y \textbf{reflexionar sobre cómo se interrelacionan} entre sí. Esta actividad tiene como propósito que, antes de profundizar en los modelos formales y las implementaciones concretas, cada alumno cuente con una visión global e integradora de los desafíos que enfrenta el diseño de redes de alta velocidad, independientemente de su escala o ámbito de aplicación.

\noindent\textbf{Fecha de entrega:} \underline{\hspace{4cm}} \quad \textbf{Valor:} \underline{\hspace{2cm}}~\% de la calificación parcial.

\section*{Instrucciones generales}
\begin{itemize}
    \item El reporte debe redactarse en formato digital (PDF), con interlineado de 1.5, fuente de 12 pt y márgenes de 2.5 cm (superior/inferior) y 3 cm (izquierdo/derecho).
    \item Extensión mínima: \textbf{4 páginas} de contenido (sin contar portada ni referencias).
    \item Se evaluará la claridad expositiva, la profundidad del análisis y la correcta identificación de conceptos.
    \item Las figuras o diagramas deben incluirse con su respectiva referencia al artículo original.
    \item Se permite y fomenta la inclusión de esquemas propios elaborados por el alumno.
\end{itemize}

\newpage

% ==========================================
% ESTRUCTURA REQUERIDA DEL REPORTE
% ==========================================
\section*{Estructura requerida del reporte}
A continuación se detalla cada sección que debe contener el reporte. El estudiante sustituirá los textos entre corchetes \texttt{[...]} por su propio contenido.

\section{Resumen de la lectura}
\textit{[Redactar un resumen de entre 200 y 300 palabras que capture la idea central del artículo: el papel de las interconexiones de alta velocidad en la era de Internet, desde el enlace punto a punto hasta las redes en chip. Mencionar el enfoque mixto ascendente/descendente empleado por el autor.]}

\section{Conceptos fundamentales identificados}
\textit{[Listar y definir con palabras propias los conceptos principales abordados en el artículo. Se sugiere incluir, como mínimo, los siguientes:]}
\begin{itemize}
    \item Ancho de banda y segmentación de canales (pipelining de enlace).
    \item Control de flujo y unidades de control de flujo (\textit{flits}).
    \item Gestión de colas y planificación de transmisión/recepción.
    \item Calidad de servicio (QoS) y control de admisión.
    \item Arquitectura de un conmutador: puertos, colas, conmutador interno, unidad de enrutamiento y arbitraje.
    \item Formato de paquete y técnicas de conmutación (conmutación de paquetes vs. \textit{wormhole}).
    \item Bloqueo de cabeza de línea (HOL blocking).
    \item Topologías: redes directas (malla), redes multietapa (\textit{fat tree}/Benes), redes irregulares.
    \item Algoritmos de enrutamiento: determinista, adaptativo, multidestino.
    \item Interbloqueo (\textit{deadlock}).
    \item Equilibrio de carga y control de congestión.
    \item Tolerancia a fallos y reconfiguración de red.
    \item Interfaz de red y sus servicios.
\end{itemize}

\section{Problemáticas de diseño al escalar una red}
\textit{[Seleccionar \textbf{al menos tres} problemáticas de diseño que surgen al aumentar el tamaño y la complejidad de una red. Para cada una, se debe:]}
\begin{enumerate}
    \item Describir la problemática con claridad.
    \item Explicar la causa raíz según lo expuesto en el artículo.
    \item Mencionar las soluciones o mecanismos propuestos.
    \item Incluir un ejemplo o esquema ilustrativo (propio o adaptado del texto).
\end{enumerate}

\textit{[Ejemplos sugeridos de problemáticas a elegir:]}
\begin{itemize}
    \item Interbloqueo (Fig.~8 del artículo).
    \item Congestión y formación de árboles de congestión.
    \item Entrega de paquetes fuera de orden por enrutamiento adaptativo.
    \item Bloqueo de cabeza de línea (HOL blocking).
    \item Degradación de la fiabilidad al crecer el número de componentes.
    \item Limitaciones de ancho de banda en un bus compartido.
\end{itemize}

\section{Interrelación entre las problemáticas identificadas}
\textit{[Analizar cómo las problemáticas elegidas en la Sección~3 se relacionan entre sí. Responder a preguntas como: ¿la solución a una problemática puede generar otra? ¿Existen mecanismos que aborden varias problemáticas simultáneamente? Se recomienda incluir un diagrama de relaciones elaborado por el alumno.]}

\textit{[Ejemplo de guía: el enrutamiento adaptativo mejora el equilibrio de carga, pero introduce entrega fuera de orden; las restricciones para evitar interbloqueo pueden limitar las rutas disponibles y agravar la congestión; el control de admisión para QoS puede mitigar la congestión, etc.]}

\section{Reflexión personal}
\textit{[Escribir una reflexión de entre 200 y 300 palabras sobre:]}
\begin{itemize}
    \item ¿Qué aspecto del artículo resultó más relevante o sorprendente?
    \item ¿Cómo se relaciona lo leído con la experiencia cotidiana del estudiante como usuario de Internet?
    \item ¿Qué temas se desean profundizar a lo largo del curso?
\end{itemize}

\section{Conclusiones}
\textit{[Sintetizar en un párrafo las ideas principales extraídas de la lectura y su importancia como base para el resto de la materia.]}

% ==========================================
% REFERENCIAS
% ==========================================
\begin{thebibliography}{9}
    \bibitem{duato2016}
    J. F. Duato Marín, ``Redes de comunicaciones de alta velocidad: ¿cómo funcionan?'', \textit{Rev. R. Acad. Cienc. Exact. Fís. Nat. (Esp)}, vol. 109, n.º 1--2, pp. 75--86, 2016.
    \bibitem{extra}
    \textit{[Cualquier referencia adicional consultada por el alumno.]}
\end{thebibliography}

\appendix
\section{Material complementario (opcional)}
\textit{[Espacio para incluir diagramas adicionales, tablas comparativas de topologías, capturas de figuras del artículo con comentarios propios, etc.]}

\newpage
% ==========================================
% RÚBRICA DE EVALUACIÓN
% ==========================================
\section*{Rúbrica de evaluación}

\begin{table}[h]
    \centering
    \begin{tabular}{@{}p{7cm}cc@{}}
        \toprule
        \textbf{Criterio} & \textbf{Puntaje} & \textbf{Obtenido} \\
        \midrule
        Resumen claro y completo de la lectura & 15\,\% & \\
        Identificación correcta de conceptos clave & 20\,\% & \\
        Análisis de $\geq 3$ problemáticas de diseño & 25\,\% & \\
        Interrelación y diagrama de relaciones & 15\,\% & \\
        Reflexión personal argumentada & 10\,\% & \\
        Redacción, formato y ortografía & 10\,\% & \\
        Puntualidad en la entrega & 5\,\% & \\
        \midrule
        \textbf{Total} & \textbf{100\,\%} & \\
        \bottomrule
    \end{tabular}
\end{table}

\vspace{1.5cm}
\noindent\rule{\textwidth}{0.4pt}
\vspace{0.5cm}

\noindent\textit{Nota: Este documento constituye la guía de entrega. El reporte final debe contener únicamente las secciones 1 a 7 y, opcionalmente, el apéndice. La portada y este enunciado no forman parte de la extensión mínima requerida.}

\end{document}
